\section{The scope of the ind-object theorem}

\begin{claim}{B.1}{Proved}{Descent and the two proofs of R3}
The specified comparison in both supplied proofs of R3 is valid. The first proof also applies to any $\U$-small colimit of representables. In the second proof, filteredness is used precisely to preserve the injections induced by an acyclic map.
\end{claim}
\begin{proof}
Here are the individual identifications, including the comparisons implicit in the supplied references. Let $q:\Cathot_\U\to\Hot_\U$ send $X$ to $BX$. A natural transformation $u\Rightarrow v$ is represented by a functor $X\times[1]\to Y$. The projection $X\times[1]\to X$ is a nerve equivalence, and its two sections become equal on localization. Thus localization of $\Cat$ at nerve equivalences factors through $\Cathot$; its universal property proves that $q$ is the corresponding localization.

For an arbitrary small diagram $K:I\to\Cathot$, put $H=\colim_i h_{K_i}$. A presheaf $H$ that inverts nerve equivalences is $q^*F$ for a unique presheaf $F$ on $\Hot$, up to its specified natural isomorphism. For any presheaf $T$ on $\Hot$, Yoneda gives
\[
 \operatorname{Nat}(\colim_i h_{BK_i},T)
 =\lim_iT(BK_i)
 =\operatorname{Nat}(H,q^*T)
 =\operatorname{Nat}(F,T).
\]
The last equality uses full faithfulness of restriction along a localization. Yoneda in the presheaf category therefore identifies
\[
 F(Y)=\colim_i[Y,BK_i]. \tag{B.1.1}
\]
Tracing the universal elements proves that the identification sends a functor $X\to K_i$ to its classifying map. This checks CATHOT A.7 and the first paragraph of GALOIS B.12, including naturality. A larger ambient universe for presheaves prevents any smallness assumption on the collection of all objects of $\Cathot_\U$.

Choose an infinite cardinal $\mu\in\U$ bounding every $|\Mor K_i|$. Fundamental groups of $BK_i$ also have cardinality at most $\mu$: their generators are edges of the nerve and their relations are triangular faces. For a small CW space $Y=\coprod_cY_c$, the checked R2 construction supplies acyclic maps $t_c:BG_c\to Y_c$ with quotient homomorphisms $p_c:G_c\twoheadrightarrow\pi_1Y_c$, such that every homomorphism from $G_c$ into a group of size at most $\mu$ kills $\ker p_c$. Put $T=\coprod_cBG_c$, $P=\coprod_cB\pi_1Y_c$, and write $v:Y\to P$ for first truncation.

For every $i$, precomposition with $p=\coprod_cp_c$ is an actual isomorphism
\[
 \Fun\bigl(\coprod_c\mathbf B\pi_1Y_c,K_i\bigr)
 \longrightarrow
 \Fun\bigl(\coprod_c\mathbf BG_c,K_i\bigr). \tag{B.1.2}
\]
A functor on the right specifies an object $k_c$ and a homomorphism into $\Aut_{K_i}(k_c)$, which factors uniquely through $p_c$. A natural transformation is a family of arrows intertwining the actions. Surjectivity of $p_c$ makes its two intertwining conditions identical. This argument proves an isomorphism before taking components, including infinitely many components of $Y$.

The first proof now uses the checked full mapping-class assertion of R2: $[Y,BK_i]\to[T,BK_i]$ is bijective. The composites $vt$ and $Bp$ agree in $\Hot$, since maps into a groupoid nerve are classified by their fundamental-groupoid functors. Taking the colimits in (B.1.1)--(B.1.2) yields a commuting triangle in which $F(t)$ and $F(vt)$ are bijections. Thus $F(v)$ is a bijection. No filteredness is used in this paragraph.

For the second proof only injectivity of $[Y,BK_i]\to[T,BK_i]$ is needed. To check its cited GALOIS B.9, form the double mapping cylinder $D=Y\cup_T^hY$. Surjectivity of $\pi_1t$ gives $\pi_1D\cong\pi_1Y$ by van Kampen. The local-coefficient Mayer--Vietoris sequence gives a homology isomorphism for the fold $D\to Y$. With coefficients $\mathbb Z[\pi_1Y]$ this is a homology equivalence on simply connected universal covers, hence a weak equivalence by homological Whitehead. The two inclusions $Y\to D$ are consequently equal in $\Hot$, proving injectivity on mapping classes. This verifies each step of B.9. For filtered $I$, equality of two elements of a colimit is witnessed at a common later stage; injectivity at that stage proves injectivity of $F(t)$. Since $F(vt)$ is bijective, cancellation proves $F(v)$ bijective. This is precisely the argument of GALOIS B.12.

Finally, for groupoids $G$, restriction gives an isomorphism $\Fun(\Pi_1X,G)\cong\Fun(X,G)$: both functors and natural transformations extend uniquely over the formal inverses. Also
\[
 N\Fun(X,G)_r=\Fun(X\times[r],G)
 =\Hom_{\sSet}(NX\times\Delta[r],NG),
\]
and $NG$ is Kan. These equalities check CATHOT C.1 and the reflection in C.4. For filtered $I$, orthogonality to all units $\eta_X:X\to\Pi_1X$ supplies, after taking the limit over $i$, a morphism $r:(\Pi_1K_i)_i\to(K_i)_i$ with $r\eta=\id$. The target ind-groupoid is also orthogonal to the units. Precomposing $\eta r$ and its identity with $\eta$ therefore proves $\eta r=\id$. This checks the specified canonical ind-isomorphism, with no change of the index category.
\end{proof}

\begin{claim}{B.2}{Proved}{All invariant small colimits of representables}
The $\Winfty$-invariant presheaves on $\Cathot_\U$ that are $\U$-small colimits of representables are exactly the presheaves
\[
 X\longmapsto\colim_{i\in I}\pi_0\Fun(\Pi_1X,G_i),
\]
where $I$ is any $\U$-small category and $G:I\to\Ho(\Grpd_\U)$ is a diagram. For an invariant presentation $H=\colim_i h_{K_i}$ the canonical comparison
\[
 \colim_i h_{K_i}\longrightarrow\colim_i h_{\Pi_1K_i} \tag{B.2.1}
\]
is an isomorphism of presheaves, for the original index $I$.
\end{claim}
\begin{proof}
The first proof checked in B.1 gives $H(\eta_X)$ bijective for every $X$. Let $L=\Pi_1:\Cathot\to\Ho(\Grpd)$ and let $j$ be inclusion. Left Kan extension along $L^{\op}$ sends $h_K$ to $h_{LK}$, by Yoneda. Its unit on $H$ is therefore (B.2.1). Since $H$ inverts every unit of $L\dashv j$, it factors as $(Hj)L^{\op}$, and the unit of this Kan extension is invertible. Conversely, each groupoid representable inverts $\Winfty$, and a colimit of bijections of sets is a bijection. The empty indexing category is included: it yields the everywhere-empty presheaf on both sides.
\end{proof}

\begin{claim}{B.3}{Proved}{Truncation localizers}
For every integer $n\geq1$, the $\W_n$-invariant $\U$-small ind-objects are exactly the ind-groupoids. The $\W_0$-invariant ones are exactly the ind-objects of discrete sets. Their canonical representatives are respectively $(\Pi_1K_i)_i$ and $(\pi_0K_i)_i$. For representable presheaves the corresponding classes consist of objects isomorphic in $\Cathot$ to groupoids and to discrete categories.
\end{claim}
\begin{proof}
Every $\W_n$ contains $\Winfty$, so B.1 first gives the canonical ind-groupoid. For $n\geq1$, a $\W_n$ map induces an equivalence on fundamental groupoids, and hence is inverted by each groupoid representable. For $n=0$, apply the same unit argument to the reflection $\pi_0:\Ho(\Grpd)\to\Set$. Its unit $G\to\pi_0G$ belongs to $\W_0$, and its target representables invert exactly the component comparisons. The same proof with a constant one-object index gives the representable assertions.
\end{proof}

\begin{claim}{B.4}{Proved}{Integral homology and homologically local groups}
Let $\mathcal G_{\mathbb Z}$ be the full subcategory of groupoids whose isotropy groups are local for the homomorphisms inducing an isomorphism on $H_1(-;\mathbb Z)$ and a surjection on $H_2(-;\mathbb Z)$. The $\W_{\mathbb Z}$-invariant $\U$-small ind-objects of $\Cathot$ are exactly the ind-objects of $\mathcal G_{\mathbb Z}$. Their canonical presentation is $(E_{\mathbb Z}\Pi_1K_i)_i$, where $E_{\mathbb Z}$ is groupwise integral homological localization. The invariant representables are exactly the categories isomorphic in $\Cathot$ to objects of $\mathcal G_{\mathbb Z}$.
\end{claim}
\begin{proof}
Bousfield's Theorems 5.2 and 5.5 and Lemma 6.1 in \citet{Bousfield1975}, opened including their proofs and applied to $R=\mathbb Z$, supply the precise ingredients: groupwise localization exists, its maps are isomorphisms on $H_1$ and surjections on $H_2$, and every such map $\pi_1X\to H$ is realized by a homology equivalence $X\to Y$. The constructions take place in $\U$. For the last ingredient, attach $1$- and $2$-cells realizing the homomorphism; the relative homology has only degree two, is free there, and the Hopf exact sequence makes its generators spherical under the $H_2$-surjectivity assumption. Attach $3$-cells on those generators. This verifies the needed realization comparison.

The reflection of groups induces a reflection on groupoids modulo natural isomorphism. Indeed a homomorphism $G\to H$ into a local group extends uniquely over $G\to E_{\mathbb Z}G$. An element conjugating two restrictions also conjugates their unique extensions, so the adjunction passes to unpointed functor classes.

For a homology equivalence of connected spaces $X\to Y$, the Hopf exact sequence shows that $\pi_1X\to\pi_1Y$ belongs to the specified group class. Consequently local-groupoid representables invert such maps, including disconnected spaces by taking componentwise functors. Conversely, B.1 first replaces an invariant ind-object by its canonical ind-groupoid. For each group $G$, realize $G\to E_{\mathbb Z}G$ by a homology equivalence $BG\to Y$ as above. Its composite with $Y\to B\pi_1Y$ is $BG\to BE_{\mathbb Z}G$. The given presheaf inverts the first map by homology invariance and the second by B.1. It therefore inverts the units of the groupoid reflection. The ind-unit argument proves the stated canonical isomorphism. A constant diagram gives the representable assertion.
\end{proof}

\begin{claim}{B.5}{Proved}{The enlarged-universe classification}
Restriction along $q:\Cathot_\U\to\Hot_\U$ gives an equivalence from $\Ind_\V(\Hot_\U)$ onto the full subcategory of $\Winfty$-invariant objects in $\Ind_\V(\Cathot_\U)$. In particular every $\U$-small homotopy type has a $\V$-small filtered ind-presentation on $\Cathot_\U$.
\end{claim}
\begin{proof}
First fix a $\U$-small homotopy type $Y$. The category of elements of $X\mapsto[BX,Y]$ on $\Cathot_\U$ is $\V$-small. It is nonempty, since it contains $(\varnothing,!)$. Disjoint union gives an upper bound for any two objects. For two parallel arrows $(A,a)\rightrightarrows(B,b)$ represented by $u,v:A\to B$, form the Grothendieck construction $C$ of the diagram with two parallel arrows $A\rightrightarrows B$. Its objects are those of $A$ and $B$; its cross arrows from $a_0\in A$ to $b_0\in B$ are pairs consisting of a choice of $u$ or $v$ and an arrow from the chosen image of $a_0$ to $b_0$. The two tagged identity arrows give natural transformations from the inclusion of $A$ to the two composites $wu,wv$, so $wu=wv$ in $\Cathot$.

The nerve homotopy-colimit comparison identifies $BC$ with the double mapping cylinder of $Bu,Bv$, compatibly with $BB$. This comparison follows directly by realizing the simplicial replacement of this finite diagram; each summand records precisely an internal arrow or one tagged cross arrow as above. The equality $bBu=bBv$ in $\Hot$ supplies a homotopy and hence an extension $c:BC\to Y$ with $cBw=b$. Thus $(B,b)\to(C,c)$ equalizes the two arrows. The element category is filtered. The density map sends a functor $X\to C$ at $(C,c)$ to $cBf$; the identity at $(X,x)$ proves surjectivity and identifies any two representatives with equal image. This checks CATHOT D.3 and D.7 for the precise maps used here.

Every object of $\Ind_\V(\Hot_\U)$ is a filtered colimit of such representables. Its restriction remains a $\V$-small filtered colimit of representables on $\Cathot_\U$: a filtered colimit of flat presheaves is flat, as witnessed directly by the common-stage and equalization conditions for its category of elements. All the indexing data are $\V$-small. Conversely, B.1.1 for a $\V$-small filtered presentation of an invariant presheaf gives its descended object in $\Ind_\V(\Hot_\U)$; this step needs no cardinal bound. Restriction is fully faithful because $q$ is a localization. The two constructions are inverse by their stated canonical comparisons.
\end{proof}
\begin{claim}{B.6}{Proved}{An exact algebraic classification for every larger localizer}
Let $\W\supseteq\Winfty$ be any fundamental localizer in $\Cat_\U$. Let $\mathcal G=\Ho(\Grpd_\U)$ and let $S_{\W}$ be the class of arrows $\Pi_1u$ for $u\in\W$. Form the ordinary localization $q_{\W}:\mathcal G\to\mathcal D_{\W}=\mathcal G[S_{\W}^{-1}]$ in $\V$. Under the fully faithful embedding of ind-groupoids as presheaves, the exact class of invariant small ind-objects is
\[
 \Ind_\U(\mathcal G)\ \cap\
 q_{\W}^*\operatorname{PSh}(\mathcal D_{\W}). \tag{B.6.1}
\]
Here the intersection is an intersection of essential images inside presheaves on $\mathcal G$. Its membership is given by the explicit extension and equality tests below. For representables, it consists of groupoids $G$ satisfying the displayed algebraic bijection (B.6.2). These descriptions require no accessibility or reflectivity assumption on $\W$.
\end{claim}
\begin{proof}
B.1 gives the canonical ind-groupoid for every invariant small ind-object. A presheaf on $\mathcal G$ descends along $q_{\W}$ exactly when it sends every arrow of $S_{\W}$ to a bijection. Reflection $\Pi_1$ then identifies this condition with invariance under $\W$, giving both directions of (B.6.1).

Here is a fully algebraic form of the criterion. Replace each stage by a skeletal groupoid with component set $T_i$ and isotropy groups $G_{i,t}$. For a groupoid $A$ put
\[
 R_i(A)=\coprod_{f:\pi_0A\to T_i}
          \prod_{a\in\pi_0A}
       \Hom(\pi_1(A,a),G_{i,f(a)})/G_{i,f(a)}. \tag{B.6.2}
\]
The quotient in each factor is by target conjugation. Chosen base objects identify this formula naturally up to its canonical conjugations; it is exactly $\pi_0\Fun(A,G_i)$, since natural isomorphisms are componentwise conjugating arrows. For $s:A\to B$ in $S_{\W}$, require these two conditions. Every element of $R_i(A)$ extends, after some transition $i\to j$, to an element of $R_j(B)$; and any two elements at a common stage of $R_i(B)$ whose restrictions become equal at a later stage themselves become equal at a further stage. For a general filtered index, ``later'' means an actual arrow, and common stages exist by filteredness. The first condition is surjectivity of $\colim_iR_i(B)\to\colim_iR_i(A)$; the second is its injectivity, by the finite-witness description of equality in filtered colimits. This proves an explicit stage criterion for (B.6.1).

For a representable, delete the index and require the restriction function $R(B)\to R(A)$ to be bijective for every $s\in S_{\W}$. This classifies all the requested representables by component sets and groups. B.3 and B.4 compute this criterion for the specified localizers. The construction of $\mathcal D_{\W}$ is legitimate because $\mathcal G$ and its arrow class are $\V$-small; no reflector in $\Grpd_\U$ is asserted for an arbitrary class.
\end{proof}

\begin{claim}{B.7}{Proved}{The sharp cardinal threshold for bounded tests}
Let $\mu,\kappa\in\U$ be infinite cardinals, and let $K:I\to\Cathot_\U$ have $\U$-small filtered index and $|\Mor K_i|\leq\mu$ for every $i$. Among such presentations, invariance on all categories with fewer than $\kappa$ arrows forces the canonical global ind-groupoid isomorphism for every $K$ exactly in the range
\[
                         \mu^+<\kappa. \tag{B.7.1}
\]
Outside this range, every small homotopy type occurs on those tests using stages with at most $\mu$ arrows. In particular $K(C_2,2)$ gives a presentation that is invariant on the prescribed tests and has no global ind-groupoid isomorphism.
\end{claim}
\begin{proof}
We first check the size-controlled localization needed for the positive assertion. For any infinite cardinal $\nu\in\U$, let $\mathcal C_{\leq\nu}$ be a skeleton of categories with at most $\nu$ arrows and let $\mathcal T_{\leq\nu}$ be homotopy types admitting simplicial models with at most $\nu$ simplices. The nerve and the category of simplices give a localization equivalence
\[
 \mathcal C_{\leq\nu}[W_\infty^{-1}]
       \simeq\mathcal T_{\leq\nu}. \tag{B.7.2}
\]
Indeed a nerve has at most $\nu$ simplices, and a category of simplices has at most $\nu$ arrows. The last-vertex comparison for a category and the realization comparison for a simplicial set are natural weak equivalences. A Kan replacement by successive finite horns has at most $\nu$ simplices, so representatives of maps and homotopies in this subcategory stay within the same bound. This proves the full localization comparison, not just essential surjectivity.

Set $\nu=\mu^+<\kappa$. The restricted presheaf on $\mathcal C_{\leq\nu}$ is invariant and has the same representable presentation, because all $K_i$ belong to that category. The descent calculation of B.1 applies there. For $Y=BK_i$, R2 has a group of cardinality at most $\max(\aleph_0,\mu^+,|NK_i|)=\nu$; its fundamental-group quotient and every intermediate nerve model also fit (B.7.2). The first proof of B.1 therefore proves $H(\Pi_1K_i)\to H(K_i)$ bijective for every stage. Taking their limit constructs the inverse to the canonical ind-unit, exactly as in B.1. The resulting isomorphism is in the global ind-category.

For the converse, suppose $\kappa\leq\mu^+$. Use the category of elements of $C\mapsto[BC,Y]$ on $\mathcal C_{\leq\mu}$. The explicit coproduct and double-mapping-cylinder Grothendieck construction of B.5 have at most $\mu$ arrows, so its proof makes this category filtered. It is $\U$-small: there is a set of category structures on sets of size at most $\mu$, and all their mapping-class sets into $Y$ are small in $\U$. Density gives an ind-presentation with these same categories as stages, representing $[B-,Y]$ on $\mathcal C_{\leq\mu}$. Every test with fewer than $\kappa$ arrows lies there. Take $Y=K(C_2,2)$. A finite face-poset model $P$ of $S^2$ gives two elements, whereas the terminal category gives one. Every ind-groupoid takes $P\to *$ to a bijection. Thus this particular bounded invariant presentation cannot be globally isomorphic to an ind-groupoid. This proves sharpness, including equality in the boundary $\kappa=\mu^+$.
\end{proof}

\begin{claim}{B.8}{Proved}{All bounded invariant restrictions}
Fix an infinite $\kappa\in\U$, and let $\mathcal C_{<\kappa}$ be a skeleton of categories with fewer than $\kappa$ arrows, with transformation-homotopy morphisms. Its localization is denoted $\mathcal D_{<\kappa}$. The restrictions of small ind-objects that are invariant on these tests are exactly the presheaves on $\mathcal D_{<\kappa}$ whose pullbacks $F$ satisfy the following finite conditions:
\begin{enumerate}[label=(\roman*)]
\item $F(\varnothing)$ is a singleton and $F(A\amalg B)\to F(A)\times F(B)$ is bijective;
\item for every pair $u,v:A\rightrightarrows B$, with $C$ the specified Grothendieck homotopy coequalizer of B.5 and $w:B\to C$, the function
\[
 F(C)\longrightarrow\Eq\bigl(F(B)\rightrightarrows F(A)\bigr)
\]
is surjective.
\end{enumerate}
Every such restriction has a presentation by stages with fewer than $\kappa$ arrows. For $\kappa>\aleph_0$, these classes are equivalently $\Ind_\U(\Hot_{<\kappa})$, where a type belongs to $\Hot_{<\kappa}$ when it has a simplicial model of size less than $\kappa$. For $\kappa=\aleph_0$, the two displayed finite conditions give the classification directly on the localization of finite categories.
\end{claim}
\begin{proof}
Let $H=\colim_i h_{K_i}$ have invariant restriction. Finite products commute with filtered colimits of sets, so (i) holds. Take $b\in H(B)$ equalizing $u,v$. At a common stage it is represented by a functor $B\to K_j$ and a finite zigzag of natural transformations between its composites with $u,v$. Encode that zigzag by a fence $Z$. Form a finite indexing category with the vertices of $Z$, one extra vertex carrying $B$, and arrows from both endpoints to that vertex. Put $A$ at each fence vertex, identity transition functors along its edges, and $u,v$ on its two end arrows. Its Grothendieck construction $C_Z$ has fewer than $\kappa$ arrows. The given transformations give an actual functor $C_Z\to K_j$ extending $b$.

Collapse the fence to the central $A$ vertex in the two-arrow diagram. This is a functor $C_Z\to C$ equal to the identity on $B$. On the nerve homotopy-colimit models it collapses a subdivided interval to the standard interval in the double mapping cylinder, so it is a weak equivalence, relative to $BB$. Invariance transfers the element of $H(C_Z)$ to $H(C)$ and proves (ii).

Conversely, (i) makes the category of elements of $F$ nonempty and gives upper bounds for pairs. For parallel arrows in that element category, (ii) gives an extension at $C$, and the actual natural transformations in B.5 equalize those arrows in $\mathcal C_{<\kappa}$. Thus the element category is filtered. It is $\U$-small, and density supplies the stated presentation by categories of size less than $\kappa$. This proves both directions for every infinite $\kappa$.

For uncountable $\kappa$, the size check in (B.7.2) works for every individual size less than $\kappa$: adjoining countably many finite pieces to an object of infinite size $\lambda$ has size at most $\max(\lambda,\aleph_0)<\kappa$. Thus $\mathcal D_{<\kappa}\simeq\Hot_{<\kappa}$. Every pair of homotopy maps has simultaneous categorical representatives within that bound. The two finite conditions are therefore precisely filteredness of the element category on $\Hot_{<\kappa}$, using finite coproducts and homotopy coequalizers. Equivalently these are its small ind-objects. This final identification does not assume that the nerve of a finite category is a finite simplicial set.
\end{proof}

